%LTeX: language=pt-BR
\section{Objetivos}

O objetivo principal deste estágio de pesquisa no exterior é estender e acelerar a resolução de problemas de otimização com regularizações descontínuas (tais como a penalidade $L_0$ e restrições de cardinalidade em grupos) por meio da exploração de novas alternativas de segunda ordem e operadores proximais estruturados. A intenção é utilizar esses métodos como poderosos pré-processadores globais capazes de navegar eficientemente por paisagens complexas e se aproximar da variedade ativa correta. Essa identificação aprimorada de boas regiões permitirá que refinadores locais --- como a descida de coordenadas e buscas combinatórias --- encontrem soluções de maior qualidade com muito mais facilidade.

Os objetivos específicos incluem:
\begin{itemize}
    \item \textbf{Desenvolvimento Algorítmico Híbrido:} Projetar um arcabouço algorítmico híbrido onde novas alternativas estruturadas de segunda ordem e operadores proximais (ex: métricas variáveis tridiagonais em blocos) atuem como identificadores globais de boas regiões. Isso será integrado à descida de coordenadas e às buscas locais combinatórias estudadas no doutorado, as quais atuarão como refinadores locais para isolar o suporte ótimo final;
    \item \textbf{Análise Teórica de Convergência:} Analisar a convergência local e global do método híbrido resultante, explorando as recentes análises de complexidade para otimização não-suave \cite{leconte2023complexity};
    \item \textbf{Estudo de Casos Práticos:} Validar os algoritmos propostos em aplicações reais de aprendizado de máquina esparso e processamento de sinais em grande escala;
    \item \textbf{Integração de Software:} Implementar os algoritmos resultantes em Julia, integrando as metodologias ao ecossistema \texttt{JuliaSmoothOptimizers} (JSO) se aplicável, garantindo alta escalabilidade e disponibilizando as ferramentas para a comunidade científica.
\end{itemize}

\subsection{Caminhos Alternativos de Pesquisa}
Embora o foco principal permaneça no arcabouço algorítmico híbrido, este projeto também identifica dois caminhos de pesquisa alternativos e promissores:
\begin{itemize}
    \item \textbf{Exploração de Separabilidade Inerente:} Aplicar o método de gradiente proximal com métrica variável \cite{Leconte_2024}, em conjunto com a descida de coordenadas e as buscas combinatórias, em problemas onde a Hessiana verdadeira é naturalmente parcialmente separável ou separável em blocos. Esta abordagem evita a necessidade de aproximações quasi-Newton que imponham tais critérios, aproveitando a estrutura nativa da função objetivo;
    \item \textbf{Técnicas Combinatórias para Otimização com Restrições:} Expandir as técnicas de busca combinatória para problemas com restrições (focando inicialmente em restrições de caixa). Neste contexto, a decisão de sair de uma face da região viável será tratada de forma análoga à decisão de adicionar ou remover parâmetros do suporte ativo em problemas $L_0$ sem restrições.
\end{itemize}

